% Paper 3 — Results (borrador 2026-10-17). Fuente de cada número:
%   docs/paper3-false-acceptance.json        (curva de reglas publicadas)
%   docs/paper3-per-config-threshold.json    (heterogeneidad, FAR fijo vs re-derivado, curva δ)
%   docs/paper3-v3-robustness.json           (estimador, H0, pares, bootstrap)
%   docs/paper3-v3-power.json                (potencia)
% Todos sobre docs/paper3-archive-manifest.txt.

\section{Results}\label{sec:results}

\subsection{Noise is a property of the configuration}\label{sec:res-heterogeneity}
Across the 149 arms of the 26 calibration sweeps, the share of tasks that flip
between two runs of the same configuration has a median of 5.7\%, a 90th
percentile of 17.6\% and a maximum of 33.3\%; 36 arms (one in four) never flip. Part of
this spread is between suites, but most of it survives inside a single sweep,
where suite, node and harness version are shared: there, the flip rates of
different configurations differ by a median of 11.1 percentage points, and by
26.3 points in the widest sweep (from 0\% to 26.3\% on the same suite). Against the
tolerance of its own sweep's base, 29 of the 90 incumbents (32\%) are noisier than
the base.

\subsection{A tolerance calibrated once}\label{sec:res-fixed}
Table~\ref{tab:fixed} reports the false-acceptance rate of the two rules.
With the primary estimator and the first arm as base, the fixed tolerance commits
6.9\% of null changes and the re-derived tolerance 0.8\%. The difference is
concentrated where the base gives a misleading picture of the noise. In 14
incumbents whose base never flipped across its repetitions, the fixed tolerance
is zero and the rule commits any positive $\Delta$: 25.7\% of their null pairs,
against 0.4\% with the re-derived tolerance. In the other 76 incumbents the rates
are 3.5\% and 0.8\%.

\begin{table}[t]
\centering
\caption{False-acceptance rate on null pairs of incumbents (1{,}814 pairs, 90
incumbents, 26 sweeps). The 95\% intervals of the difference come from a cluster
bootstrap over sweeps.}\label{tab:fixed}
\begin{tabular}{llrrr}
\toprule
$\hat\delta$ estimator & base $H_0$ & fixed & re-derived & difference [95\% CI] \\
\midrule
max $|\Delta|$ & first arm & 6.9\% & 0.8\% & $+6.1$ [$+2.2$, $+11.3$] \\
max $|\Delta|$ & random (200 draws) & 7.3\% & 0.7\% & $+6.6$ [$+5.2$, $+8.1$] \\
max $|\Delta|$ & quietest arm & 17.8\% & 0.9\% & $+16.9$ [$+11.7$, $+22.7$] \\
max $|\Delta|$ & noisiest arm & 0.0\% & 0.5\% & $-0.5$ [$-0.9$, $-0.2$] \\
$2\,\mathrm{sd}(\Delta)$ & random & 7.5\% & 0.8\% & $+6.8$ [$+5.4$, $+8.1$] \\
mean $|\Delta|$ & random & 17.9\% & 17.8\% & $+0.1$ [$-1.1$, $+1.1$] \\
\bottomrule
\end{tabular}
\end{table}

Three robustness checks bear on this result. First, drawing the base at random
gives the same answer as taking the first arm (7.3\% against 0.7\%), so the
result does not depend on which configuration happens to be first. Second,
subsampling the base's pairs to the number available to the incumbent changes no
rate by more than 0.3 points, so the per-configuration rule does not win by using
fewer pairs. Third, replacing the maximum with twice the standard deviation,
another tail estimator, reproduces the result. The rows for the quietest and
noisiest bases bracket the bet that a single calibration makes: it is safe only
when the base turns out to be the noisiest configuration of the evolution, and
in that case the fixed rule is conservative and pays in power
(Section~\ref{sec:res-power}). When the base is the quietest, the fixed rule
commits 17.8\% of null changes.

The last row of Table~\ref{tab:fixed} shows the limit of re-derivation. With a
central estimator (the mean of $|\Delta|$), both rules commit about 18\% of null
changes. Re-deriving a threshold does not rescue a threshold set at the wrong
level: it helps only when the threshold is a tail of the null distribution.

\subsection{Power}\label{sec:res-power}
Table~\ref{tab:power} adds true improvements of known size to the null pairs.
With the base drawn at random, the re-derived rule commits more real improvements
than the fixed rule at every effect size up to 15\% of the suite, and the two tie
at 20\%; with the first arm as base the gap is wider (at a 10\% improvement, 64\%
against 38\%). The per-configuration rule therefore dominates on both axes:
fewer false acceptances and equal or greater power. The reason is that a single
tolerance is miscalibrated in both directions at once: too narrow for incumbents
noisier than the base, which produces the false acceptances, and too wide for
incumbents quieter than the base, which loses real improvements.

\begin{table}[t]
\centering
\caption{Share of committed changes when a true improvement $e$ (fraction of the
suite) is added to every null pair of the incumbents. $e = 0$ is the
false-acceptance rate.}\label{tab:power}
\begin{tabular}{rrrrr}
\toprule
 & \multicolumn{2}{c}{$H_0$ = first arm} & \multicolumn{2}{c}{$H_0$ random} \\
$e$ & fixed & re-derived & fixed & re-derived \\
\midrule
0.00 & 7\% & 1\% & 7\% & 1\% \\
0.03 & 20\% & 38\% & 33\% & 38\% \\
0.05 & 22\% & 39\% & 34\% & 39\% \\
0.10 & 38\% & 64\% & 56\% & 62\% \\
0.15 & 60\% & 83\% & 76\% & 81\% \\
0.20 & 73\% & 89\% & 87\% & 87\% \\
\bottomrule
\end{tabular}
\end{table}

The cost of re-derivation is runs. The per-configuration rule needs at least
three repetitions of every incumbent, where a fixed tolerance needs them only for
the base. For an evolution of $T$ rounds with $k$ repetitions, that adds $(k-1)T$
evaluations of the suite.

\subsection{The threshold curve}\label{sec:res-curve}
The decision rules of published optimizers span two orders of magnitude on the
1{,}614 null comparisons of the full archive (Figure~\ref{fig:curve}). Committing
the best score accepts 26.3\% of null changes; requiring a training gain without
a development regression, or a one-task margin with an anchor set, accepts
15.6\%; raising the margin to three tasks accepts 1.1\%; the exact paired test
accepts 0.1\%. The margin, not the presence of a rule, does the work: the same
margin-plus-anchor rule falls fourteen-fold when its margin moves from one task to
three. On the pooled null pairs of the incumbents, a tolerance of 5\% of the suite
accepts 25.1\% of null changes, 10\% accepts 9.0\%, 15\% accepts 1.9\% and 20\%
accepts 0.8\%.

% Figura pendiente: curva δ → FAR (docs/paper3-per-config-threshold.json,
% campo curve_pooled_delta_vs_far) con las cinco reglas superpuestas como puntos.
\begin{figure}[t]
\centering
\fbox{\parbox{0.9\linewidth}{\centering\vspace{3em}Figure pending: threshold
$\delta$ versus false-acceptance rate, with the five decision rules as
points.\vspace{3em}}}
\caption{Threshold curve on null comparisons.}\label{fig:curve}
\end{figure}

These rates are a property of this regime and archive. They are not estimates
of the false-acceptance rate of the systems that introduced the rules, which run
other models on other benchmarks. They do show that the margin alone moves a gate
across two orders of magnitude of error, and that no published gate reports where
on this curve it operates.
